% GENERATED FILE. Edit canonical lessons, then run npm run generate:books.
% canonical-source-hash: fnv1a64:ac3808842c3ec95f
% canonical-lessons: AR-C54-ju, AR-C54-said, AR-C54-hazin, AR-C54-tabib, AR-C54-jadd

\chapter{Hunger, Happy, Sad, Doctor, Grandfather}
\label{ch:ar-hunger-happy-sad-doctor-grandfather}
\hlchaptermodality{54}

\begin{chapteropening}
\textbf{By the end of this chapter:} \emph{I can say hunger, happy, sad, a doctor and a grandfather.}

Complete the last lesson of chapter 54: I can say hunger, happy, sad, a doctor and a grandfather.
\end{chapteropening}

\section[\texorpdfstring{\ar{جوع} (jūʿ)}{جوع (jūʿ)}]{\ar{جوع} (jūʿ) — hunger}

\label{lesson:AR-C54-ju}



\begin{quote}
Before the new one: say the Arabic for pain, then the Arabic for medicine.
\end{quote}

\subsection*{You'll want to know: \ar{جوع}}
\textbf{\ar{جوع}} — \emph{jūʿ} — \textquotedblleft{}hunger\textquotedblright{}.

\begin{cousinweb}
\textbf{\ar{جوع}} (\emph{jūʿ}) is built on the native Arabic root \emph{j-w-ʿ}, \textquotedblleft{}to be hungry\textquotedblright{} (Lane). \emph{Jūʿ} is the verbal noun, \textquotedblleft{}hunger\textquotedblright{}.
\end{cousinweb}

\begin{grammarlens}[title={What you've built}]
One more word for this chapter.
\end{grammarlens}

\subsection*{Your turn}
\begin{itemize}
\raggedright
  \item \emph{Say it:} \emph{jūʿ}
  \item \emph{Say it:} \emph{jūʿ}, once more
  \item \emph{Say it:} say \emph{jūʿ}
  \item \emph{Recall:} say the Arabic for pain, then the Arabic for medicine, then say \emph{jūʿ} again
\end{itemize}

\begin{tcolorbox}[breakable,colback=teal!4,colframe=teal!35!black,title={Before you move on}]
What does \ar{جوع} mean? (\textquotedblleft{}Hunger\textquotedblright{}.) Say it once more.
\end{tcolorbox}
\section[\texorpdfstring{\ar{سعيد} (saʿīd)}{سعيد (saʿīd)}]{\ar{سعيد} (saʿīd) — happy}

\label{lesson:AR-C54-said}



\begin{quote}
Before the new one: say the Arabic for medicine, then the Arabic for hunger.
\end{quote}

\subsection*{You'll want to know: \ar{سعيد}}
\textbf{\ar{سعيد}} — \emph{saʿīd} — \textquotedblleft{}happy\textquotedblright{}.

\begin{cousinweb}
\textbf{\ar{سعيد}} (\emph{saʿīd}) is built on the native Arabic root \emph{s-ʿ-d}, \textquotedblleft{}to be fortunate, to be happy\textquotedblright{} (Lane). \emph{Saʿīd} is the \emph{faʿīl} adjective, \textquotedblleft{}fortunate, happy\textquotedblright{}.
\end{cousinweb}

\begin{grammarlens}[title={What you've built}]
One more word for this chapter.
\end{grammarlens}

\subsection*{Your turn}
\begin{itemize}
\raggedright
  \item \emph{Say it:} \emph{saʿīd}
  \item \emph{Say it:} \emph{saʿīd}, once more
  \item \emph{Say it:} say \emph{saʿīd}
  \item \emph{Recall:} say the Arabic for medicine, then the Arabic for hunger, then say \emph{saʿīd} again
\end{itemize}

\begin{tcolorbox}[breakable,colback=teal!4,colframe=teal!35!black,title={Before you move on}]
What does \ar{سعيد} mean? (\textquotedblleft{}Happy\textquotedblright{}.) Say it once more.
\end{tcolorbox}
\section[\texorpdfstring{\ar{حزين} (ḥazīn)}{حزين (ḥazīn)}]{\ar{حزين} (ḥazīn) — sad}

\label{lesson:AR-C54-hazin}



\begin{quote}
Before the new one: say the Arabic for hunger, then the Arabic for happy.
\end{quote}

\subsection*{You'll want to know: \ar{حزين}}
\textbf{\ar{حزين}} — \emph{ḥazīn} — \textquotedblleft{}sad\textquotedblright{}.

\begin{cousinweb}
\textbf{\ar{حزين}} (\emph{ḥazīn}) is built on the native Arabic root \emph{ḥ-z-n}, \textquotedblleft{}to grieve, to be sad\textquotedblright{} (Lane). \emph{Ḥazīn} is the \emph{faʿīl} adjective, \textquotedblleft{}grieving\textquotedblright{}.
\end{cousinweb}

\begin{grammarlens}[title={What you've built}]
One more word for this chapter.
\end{grammarlens}

\subsection*{Your turn}
\begin{itemize}
\raggedright
  \item \emph{Say it:} \emph{ḥazīn}
  \item \emph{Say it:} \emph{ḥazīn}, once more
  \item \emph{Say it:} say \emph{ḥazīn}
  \item \emph{Recall:} say the Arabic for hunger, then the Arabic for happy, then say \emph{ḥazīn} again
\end{itemize}

\begin{tcolorbox}[breakable,colback=teal!4,colframe=teal!35!black,title={Before you move on}]
What does \ar{حزين} mean? (\textquotedblleft{}Sad\textquotedblright{}.) Say it once more.
\end{tcolorbox}
\section[\texorpdfstring{\ar{طبيب} (ṭabīb)}{طبيب (ṭabīb)}]{\ar{طبيب} (ṭabīb) — a doctor}

\label{lesson:AR-C54-tabib}



\begin{quote}
Before the new one: say the Arabic for happy, then the Arabic for sad.
\end{quote}

\subsection*{You'll want to know: \ar{طبيب}}
\textbf{\ar{طبيب}} — \emph{ṭabīb} — \textquotedblleft{}a doctor\textquotedblright{}.

\begin{cousinweb}
\textbf{\ar{طبيب}} (\emph{ṭabīb}) is built on the native Arabic root \emph{ṭ-b-b}, \textquotedblleft{}to treat with medicine\textquotedblright{} (Lane). \emph{Ṭabīb} is the one who treats, \textquotedblleft{}a physician\textquotedblright{}.
\end{cousinweb}

\begin{grammarlens}[title={What you've built}]
One more word for this chapter.
\end{grammarlens}

\subsection*{Your turn}
\begin{itemize}
\raggedright
  \item \emph{Say it:} \emph{ṭabīb}
  \item \emph{Say it:} \emph{ṭabīb}, once more
  \item \emph{Say it:} say \emph{ṭabīb}
  \item \emph{Recall:} say the Arabic for happy, then the Arabic for sad, then say \emph{ṭabīb} again
\end{itemize}

\begin{tcolorbox}[breakable,colback=teal!4,colframe=teal!35!black,title={Before you move on}]
What does \ar{طبيب} mean? (\textquotedblleft{}A doctor\textquotedblright{}.) Say it once more.
\end{tcolorbox}
\section[\texorpdfstring{\ar{جد} (jadd)}{جد (jadd)}]{\ar{جد} (jadd) — a grandfather}

\label{lesson:AR-C54-jadd}



\begin{quote}
Before the new one: say the Arabic for sad, then the Arabic for a doctor.
\end{quote}

\subsection*{You'll want to know: \ar{جد}}
\textbf{\ar{جد}} — \emph{jadd} — \textquotedblleft{}a grandfather\textquotedblright{}.

\begin{cousinweb}
\textbf{\ar{جد}} (\emph{jadd}) is built on the native Arabic root \emph{j-d-d}, \textquotedblleft{}a grandfather; fortune, good luck\textquotedblright{} (Lane). Lane gives \emph{jadd} both as the grandfather and as fortune.
\end{cousinweb}

\begin{grammarlens}[title={What you've built}]
One more word for this chapter.
\end{grammarlens}

\subsection*{Your turn}
\begin{itemize}
\raggedright
  \item \emph{Say it:} \emph{jadd}
  \item \emph{Say it:} \emph{jadd}, once more
  \item \emph{Say it:} say \emph{jadd}
  \item \emph{Recall:} say the Arabic for sad, then the Arabic for a doctor, then say \emph{jadd} again
\end{itemize}

\begin{tcolorbox}[breakable,colback=teal!4,colframe=teal!35!black,title={Before you move on}]
What does \ar{جد} mean? (\textquotedblleft{}A grandfather\textquotedblright{}.) Say it once more.
\end{tcolorbox}
